% BOOK-FILE-ID: f-45a181b08d92
\section{演習を別の視点から}\label{節：演習を別の視点から}
\bookproblemref{practice-entrance}{9}{発展問9},
\bookproblemref{practice-entrance}{4}{問4},
\bookproblemref{practice-entrance}{3}{問3}で使った操作を,不変量・回転・母関数という別の視点から見直す.
演習で解いた問題の裏にも,本編と同じ数学が隠れていたことが分かるだろう.

\subsection{整数方程式と不変量}
発展問9では,$u_0=3$, $u_1=3$, $u_{n+2}=3u_{n+1}-u_n$から,
$(3,u_n,u_{n+1})$が$x^2+y^2+z^2=xyz$の解になることを示した.
$n=0$では
\[
  u_n^2+u_{n+1}^2+9=3u_nu_{n+1}
\]
が成り立つ.一般の$x,y$に対して
\[
  y^2+(3y-x)^2-3y(3y-x)=x^2+y^2-3xy
\]
だから,$u_n^2+u_{n+1}^2-3u_nu_{n+1}$は漸化式を用いる前後で変わらず$-9$に保たれる.
$u_{n+2}=3u_{n+1}-u_n$を使えば,$u_nu_{n+2}=3u_nu_{n+1}-u_n^2=u_{n+1}^2+9$となり,
\[
  u_{n+2}=\frac{u_{n+1}^2+9}{u_n}
\]
とも書ける.正値性は$u_0=u_1=3$と,
$u_2=6$以降の増加から確認できる.
\sectionref*{節：変わらない量を使う}の定数$1$を$9$へ変えた漸化式だった.
元の問題は\bookterm{linear}{線形}の漸化式から保存を示し,本編の例は保存から線形の漸化式を取り出した.

\begin{bookpoint}
  \textbf{POINT}\par
  隣り合う項の組に対して,更新しても変わらない式を探し,初期値でその値を決める.
  保存される関係を使うと,非線形の式から線形の漸化式を取り出せる場合がある.
\end{bookpoint}

\subsection{中点・回転・面積}
発展問4では,平行四辺形の各辺上で同じ比率$t$ $(0<t<1)$の分点を取り,
新しい平行四辺形を作った.連続する辺ベクトルを$u,v$とする.
頂点についての式を引き算すれば,
\[
  u'=(1-t)u+tv,\qquad v'=-tu+(1-t)v
\]
である.

\paragraphbf{正方形の中点を結ぶ操作}
正方形で,向きを反時計回りに取った辺を複素数として$v=iu$と書けるときは,
\[
  u'=((1-t)+it)u
\]
である.$t=\frac12$なら拡大率は$\frac1{\sqrt2}$,回転角は$\frac\pi4$になる.
これが問題で中点を結んだときの辺長と向きに一致する.

\par\medskip\noindent\textbf{別解：平行四辺形を行列で見る}\par\nobreak\smallskip
正方形という条件を外し,最初に求めた$u',v'$の式へ戻る.
辺の組を混ぜる\bookterm{matrix}{行列}は
\[
  M=\begin{pmatrix}1-t&t\\-t&1-t\end{pmatrix}
\]
となる.\bookterm{matrix}{行列}式は$\rho=(1-t)^2+t^2$である.
$u=(u_1,u_2)$, $v=(v_1,v_2)$と書くと,底辺と高さから平行四辺形の面積は$|u_1v_2-u_2v_1|$となる.$u',v'$の式を代入すると
\BookMathLayout{\[
    u'_1v'_2-u'_2v'_1
    =\bigl((1-t)^2+t^2\bigr)(u_1v_2-u_2v_1).
  \]}{\[
    \begin{aligned}
      u'_1v'_2-u'_2v'_1\\
      {}=\bigl((1-t)^2+t^2\bigr)(u_1v_2-u_2v_1).
    \end{aligned}
  \]}
$\rho>0$なので,絶対値を取った面積も一回ごとに$\rho$倍となる.
この行列式の値は,\bookterm{eigenvalue}{固有値}$1-t\pm it$の絶対値の二乗でもある.

\begin{bookpoint}
  \textbf{POINT}\quad \bookterm{matrix}{行列}が作用する対象\par
  辺の組を混ぜる\bookterm{matrix}{行列}を,平面上の相似変換(長さを一様に拡大・縮小し,角度を保つ変換)と読むには条件が要る.
  面積が何倍になるかは一般の平行四辺形でも分かるが,上の回転の説明には直交・等長の辺を使った.
\end{bookpoint}

\subsection{一定になる畳み込み和}
発展問3の$c_0=1$, $c_{n+1}=\frac{2n+1}{2n+2}c_n$から,
\[
  c_n=\frac1{4^n}\binom{2n}{n}
\]
を得た.その\bookterm{generating}{母関数}を$C(x)=\sum_{n\ge0}c_nx^n$とする.
$C(x)^2$の$x^n$の係数が,求める和$\sum_{k=0}^n c_kc_{n-k}$である.
ここでは積の微分法則と,高校範囲を超える形式的な母関数の微分を使う.初めて読むときは飛ばしてもよい.
形式的な微分は$x^n\mapsto nx^{n-1}$と定め,関数としての収束は仮定しない.
元の\bookterm{recurrence}{漸化式}を$(n+1)c_{n+1}=(n+\frac12)c_n$と書けば,
\[
  (1-x)C'(x)=\frac12C(x)
\]
を得る.積の係数は有限和なので,係数を比較することで通常の積の微分法則も使える.
$D(x)=C(x)^2=\sum_{n\ge0}d_nx^n$と置くと,
\[
  (1-x)D'(x)=2C(x)(1-x)C'(x)=D(x).
\]
$x^n$の係数を比較すれば$(n+1)d_{n+1}-nd_n=d_n$,すなわち$d_{n+1}=d_n$である.
$d_0=c_0^2=1$なので,
\[
  C(x)^2=\frac1{1-x}=\sum_{n\ge0}x^n
\]
の$x^n$の係数から
\[
  \sum_{k=0}^n c_kc_{n-k}=1
\]
を得る.有限和の変形と母関数の積は,同じ\bookterm{convolution}{畳み込み}の二つの扱い方である.

\begin{bookpoint}
  \textbf{POINT}\quad 別解の使い分け\par
  別の表現は,和・面積・不変量の意味を見せる.元の条件と答えへ戻って対応を確かめる.
\end{bookpoint}
